% Un petit graphe signé, dessiné avec la convention du manuscrit
% (PartIII/ChapII, fig. « Interprétation d'un poids w réel ») et de la
% présentation NIM du 25 juin 2026 :
%   trait plein    = interaction attractive  (W_e > 0)
%   pointillés     = interaction répulsive   (W_e < 0)
%   trait gras     = arête satisfaite par la configuration sigma
%
% Deux lectures ajoutées pour la soutenance, d'après le chapitre 11 du
% manuscrit (« Un cadre bayésien pour la détection de communautés ») :
%   - les six sommets V = {1,...,6} portent leur numéro (petit chiffre gris,
%     posé vers l'extérieur du nuage) et la ligne du haut donne la
%     configuration cachée Sigma = (+,+,-,+,-,-), chaque signe surmontant
%     le numéro du sommet auquel il appartient ;
%   - le repère (X_1,X_2) rappelle que l'attribut X d'un sommet EST sa
%     position dans le plan (modèle géométrique de Sankararaman--Baccelli,
%     § 11.3 du manuscrit, où la loi d'une arête ne dépend que de sa
%     longueur r_e = || x_i - x_j ||).
% Les positions ont été très légèrement retouchées (seuls c et f bougent) pour
% que les sept paires reliées soient exactement les sept plus courtes du nuage
% (arêtes de 1,13 à 1,35 ; plus courte paire non reliée : 1,51) : avec un
% repère affiché, le dessin se lit alors comme un vrai graphe géométrique.
    \begin{tikzpicture}[x=1cm,y=1cm,scale=0.78,every node/.style={scale=0.88},
        num/.style={font=\footnotesize,text=gris_fonce_inria,inner sep=1.6pt},
        axe/.style={gris_fonce_inria,line width=0.6pt}]

      % --- la configuration cachée, composante par composante ---
      \node[anchor=south,inner sep=0pt] at (1.33,2.08)
        {$\Sigma=\bigl(
            \underset{\textcolor{gris_fonce_inria}{1}}{+},
            \underset{\textcolor{gris_fonce_inria}{2}}{+},
            \underset{\textcolor{gris_fonce_inria}{3}}{-},
            \underset{\textcolor{gris_fonce_inria}{4}}{+},
            \underset{\textcolor{gris_fonce_inria}{5}}{-},
            \underset{\textcolor{gris_fonce_inria}{6}}{-}\bigr)$};

      % --- le repère : l'attribut X d'un sommet est sa position ---
      \draw[axe,-{Latex[length=1.6mm]}] (-0.62,-0.46) -- (3.28,-0.46)
        node[num,anchor=west,inner sep=2pt] {$X_1$};
      \draw[axe,-{Latex[length=1.6mm]}] (-0.62,-0.46) -- (-0.62,1.58)
        node[num,anchor=south,inner sep=2pt] {$X_2$};

      % --- les six sommets : position = attribut X, signe = composante de Sigma ---
      \node[spin] (a) at (0.03,1.16) {$+$};
      \node[spin] (b) at (1.31,1.55) {$+$};
      \node[spin] (c) at (2.66,1.43) {$-$};
      \node[spin] (d) at (0.30,0.05) {$+$};
      \node[spin] (e) at (1.60,0.35) {$-$};
      \node[spin] (f) at (2.69,0.08) {$-$};

      % numéros posés vers l'extérieur du nuage
      \node[num,anchor=east]       at (a.180) {1};
      \node[num,anchor=south west] at (b.55)  {2};
      \node[num,anchor=south west] at (c.40)  {3};
      \node[num,anchor=north east] at (d.215) {4};
      \node[num,anchor=north]      at (e.270) {5};
      \node[num,anchor=north west] at (f.-35) {6};

      % --- arêtes satisfaites (en gras) ---
      \draw[posedge,freeze] (a)--(b);   % attractive, (+,+)
      \draw[posedge,freeze] (a)--(d);   % attractive, (+,+)
      \draw[posedge,freeze] (e)--(f);   % attractive, (-,-)
      \draw[negedge,freeze] (b)--(c);   % répulsive,  (+,-)
      \draw[negedge,freeze] (d)--(e);   % répulsive,  (+,-)

      % --- arêtes non satisfaites (en trait fin) ---
      \draw[posedge,line width=0.6pt] (b)--(e);   % attractive, (+,-)
      \draw[negedge,line width=0.6pt] (c)--(f);   % répulsive,  (-,-)
    \end{tikzpicture}
